\begin{titlepage} 
\title{{ \large \textbf{Debt Overhang During Policy Uncertainty: \\  The Case of Gold Clauses in the 1930s}\!\!}
\thanks{We are grateful to Harjoat Bhamra, Patrick Bolton, Briana Chang, Sudheer Chava, Carola Frydman, Harry DeAngelo, Gustavo Grullon, Larry Harris, Zhiguo He, Jerry Hoberg, Ayse Imrohoroglu, Selo Imrohoroglu, Marcin Kacperczyk, Raymond Kim, Arthur Korteweg, Peter Koudijs, Zack Liu, William Mann, Kevin Murphy, Dimitris Papanikolaou, Lubos Pastor, Vincenzo Quadrini, Mark Ready, Rodney Ramcharan, Tarun Ramadorai, Alejandro Rivera, Lukas Schmid, Ivan Shaliastovich, Selale Tuzel, Jessica Wachter, Jessie Wang, and Miao Ben Zhang; seminar participants at the UBC Winter Finance Conference, WFA, AEA, the University of Kentucky Finance Conference, the Tel Aviv University Finance Conference, the SFS Cavalcade, the 14th California Corporate Finance Conference, MFA, Imperial College London, McGill University, Syracuse University, USC Marshall, and the University of Wisconsin-Madison for helpful comments; and Simon Oh, Imran Majeed, and Ally Rilling for excellent research assistance. Financial support for this research was provided by the University of Wisconsin-Madison Office of the Vice Chancellor for Research and Graduate Education with funding from the Wisconsin Alumni Research Foundation as well as Jacobs Levy Equity Management Center for Quantitative Financial Research. This paper was previously titled ``Leverage Risk and Investment: The Case of Gold Clauses in the 1930s.''}\vspace*{4mm}}

\author{Joao Gomes 
\footnote{The Wharton School, University of Pennsylvania. E-mail: gomesj@wharton.upenn.edu.} \\
\and 
Mete Kilic 
\footnote{Marshall School of Business, University of Southern California. Email: mkilic@marshall.usc.edu.}
\and 
S\'ebastien Plante
\footnote{Wisconsin School of Business, University of Wisconsin-Madison. E-mail: splante@wisc.edu.} \vspace*{10mm}}


\makeatletter
\let\Hy@WarnOptionDisabled\@gobble
\makeatother
\hypersetup{pageanchor=false}
\maketitle
\thispagestyle{empty}
\hypersetup{pageanchor=true}

\vspace*{-1cm}
\begin{abstract}
\small 
\onehalfspacing
\noindent 
We study the impact of the 1933 abrogation of gold clauses on the slow recovery of corporate investment following the Great Depression. Constitutional challenges to the abrogation exposed many firms to a potential 69\% increase in bond payments. Public firms with higher exposure to this risk reduced their investment in 1933 and 1934 followed by a recovery upon the Supreme Court's 1935 decision to uphold the abrogation. In the cross section, decreases in investment during 1933 and 1934 coincide with increases in equity payouts. Higher leverage uncertainty accounts for around one third of public firms' divestment in 1933 and 1934.
\end{abstract}
\end{titlepage}

\hypersetup{pageanchor=true}
\clearpage

\setcounter{page}{1}